\section{Test 2: Operacions, potències, arrels i proporcions}

\iniciTest{tB}{c,b,a,d,c,b,a,d,c,a,d,b}

\begin{enumerate}
\pregunta El resultat de
\[
	\frac{2}{3}+\frac{5}{6}-\frac{1}{4}
\]
és:
\begin{enumerate}
\opcio{a}{\(-\frac{5}{4}\).}
\opcio{b}{\(\frac{7}{4}\).}
\opcio{c}{\(\frac{5}{4}\).}
\opcio{d}{\(\frac{1}{4}\).}
\end{enumerate}
\feedback

\pregunta El resultat de
\[
	\left(\frac{3}{4}-\frac{2}{3}\right)\cdot\frac{6}{5}+\frac{1}{2}
\]
és:
\begin{enumerate}
\opcio{a}{\(\frac{1}{5}\).}
\opcio{b}{\(\frac{3}{5}\).}
\opcio{c}{\(\frac{2}{5}\).}
\opcio{d}{\(\frac{5}{3}\).}
\end{enumerate}
\feedback

\pregunta El resultat de
\[
	\left(\frac{5}{6}-\frac{1}{3}\right):\frac{2}{9}-\frac{3}{4}
\]
és:
\begin{enumerate}
\opcio{a}{\(\frac{3}{2}\).}
\opcio{b}{\(\frac{9}{4}\).}
\opcio{c}{\(1\).}
\opcio{d}{\(\frac{7}{4}\).}
\end{enumerate}
\feedback

\pregunta El resultat de
\[
	\frac{3}{5}
	-
	\frac{7}{10}\cdot\frac{2}{3}
	+
	\frac{1}{4}:\frac{5}{8}
\]
és:
\begin{enumerate}
\opcio{a}{\(\frac{1}{15}\).}
\opcio{b}{\(\frac{2}{5}\).}
\opcio{c}{\(\frac{7}{15}\).}
\opcio{d}{\(\frac{8}{15}\).}
\end{enumerate}
\feedback

\pregunta El resultat de
\[
	\left(-\frac{2}{3}\right)^2
	-
	\frac{5}{6}\left(\frac{3}{5}-\frac{1}{10}\right)
\]
és:
\begin{enumerate}
\opcio{a}{\(\frac{1}{12}\).}
\opcio{b}{\(-\frac{1}{36}\).}
\opcio{c}{\(\frac{1}{36}\).}
\opcio{d}{\(\frac{5}{36}\).}
\end{enumerate}
\feedback

\pregunta El resultat de
\[
	\sqrt{\frac{25}{36}}
	+
	\sqrt[3]{-\frac{8}{27}}
	-
	\left(-\frac{1}{2}\right)^3
\]
és:
\begin{enumerate}
\opcio{a}{\(-\frac{7}{24}\).}
\opcio{b}{\(\frac{7}{24}\).}
\opcio{c}{\(\frac{13}{24}\).}
\opcio{d}{\(\frac{1}{24}\).}
\end{enumerate}
\feedback

\pregunta Per a qualsevol \(x\in\Q\), es compleix:
\begin{enumerate}
\opcio{a}{\(\left(x^2\right)^3=x^6\).}
\opcio{b}{\(\left(x^2\right)^3=x^5\).}
\opcio{c}{\(\left(x^2\right)^3=x^8\).}
\opcio{d}{\(\left(x^2\right)^3=6x\).}
\end{enumerate}
\feedback

\pregunta Quina igualtat és falsa en general?
\begin{enumerate}
\opcio{a}{\(x^m\cdot x^n=x^{m+n}\).}
\opcio{b}{\((xy)^n=x^ny^n\).}
\opcio{c}{\(\left(x^m\right)^n=x^{mn}\).}
\opcio{d}{\((x+y)^2=x^2+y^2\).}
\end{enumerate}
\feedback

\pregunta La fracció generatriu de \(1{,}4\overline{6}\) és:
\begin{enumerate}
\opcio{a}{\(\frac{14}{9}\).}
\opcio{b}{\(\frac{73}{50}\).}
\opcio{c}{\(\frac{22}{15}\).}
\opcio{d}{\(\frac{17}{12}\).}
\end{enumerate}
\feedback

\pregunta La tercera proporcional de \(2\) i \(6\) és:
\begin{enumerate}
\opcio{a}{\(18\).}
\opcio{b}{\(12\).}
\opcio{c}{\(9\).}
\opcio{d}{\(3\).}
\end{enumerate}
\feedback

\pregunta La quarta proporcional de \(3\), \(5\) i \(9\) és:
\begin{enumerate}
\opcio{a}{\(9\).}
\opcio{b}{\(12\).}
\opcio{c}{\(13\).}
\opcio{d}{\(15\).}
\end{enumerate}
\feedback

\pregunta La mitjana proporcional positiva de \(4\) i \(9\) és:
\begin{enumerate}
\opcio{a}{\(5\).}
\opcio{b}{\(6\).}
\opcio{c}{\(7\).}
\opcio{d}{\(\frac{9}{4}\).}
\end{enumerate}
\feedback
\end{enumerate}
